{
\captionsetup{type=figure}
\captionof{figure}{Trajectory for instance \texttt{psf\_\_requests-2317}.}
\begin{issuebox}[psf\_\_requests-2317]We're currently solving the following issue within our repository. Here's the issue text:

ISSUE:

method = builtin\_str(method) problem

In requests/sessions.py is a command:



method = builtin\_str(method)

Converts method from

b’GET’

to

"b'GET’"



Which is the literal string, no longer a binary string.  When requests tries to use the method "b'GET’”, it gets a 404 Not Found response.



I am using python3.4 and python-neutronclient (2.3.9) with requests (2.4.3).  neutronclient is broken because it uses this "args = utils.safe\_encode\_list(args)" command which converts all the values to binary string, including method.



I'm not sure if this is a bug with neutronclient or a bug with requests, but I'm starting here.  Seems if requests handled the method value being a binary string, we wouldn't have any problem.



Also, I tried in python2.6 and this bug doesn't exist there. Some difference between 2.6 and 3.4 makes this not work right.











INSTRUCTIONS: (31 lines)

\begin{CodeVerbatim}[breaklines=true, breakanywhere=true]
(Open file: n/a)
(Current directory: /psf__requests)
bash-\$
\end{CodeVerbatim}
\end{issuebox}
\begin{agentbox}[psf\_\_requests-2317]
To address the issue, we first need to replicate the behavior described. The issue mentions that the `builtin\_str` function in `requests/sessions.py` is converting a binary string method like `b'GET'` into the literal string `''b'GET'''`, which is incorrect. We should start by locating the `sessions.py` file within the `requests` library to examine the code around the `builtin\_str` function call.



Let's begin by finding the `sessions.py` file in the current repository.

\hrulefill
\begin{CodeVerbatim}[breaklines=true, breakanywhere=true]
find_file "sessions.py"
\end{CodeVerbatim}
\end{agentbox}
\begin{observationbox}[psf\_\_requests-2317]\begin{Verbatim}[breaklines=true, breakanywhere=true]
Found 2 matches for "sessions.py" in /psf__requests:
/psf__requests/build/lib/requests/sessions.py
/psf__requests/requests/sessions.py

(Open file: n/a)
(Current directory: /psf__requests)
bash-$
\end{Verbatim}
\end{observationbox}
\begin{agentbox}[psf\_\_requests-2317]
We have found two matches for `sessions.py`. The one in the `build/lib/requests` directory is likely a build artifact, and the one we're interested in is probably `/psf\_\_requests/requests/sessions.py`. Let's open this file to inspect the code around the `builtin\_str` function call.



\hrulefill
\begin{CodeVerbatim}[breaklines=true, breakanywhere=true]
open requests/sessions.py
\end{CodeVerbatim}
\end{agentbox}
\begin{observationbox}[psf\_\_requests-2317]\begin{Verbatim}[breaklines=true, breakanywhere=true]
[File: /psf__requests/requests/sessions.py (671 lines total)]
1:# -*- coding: utf-8 -*-
2:
3:"""
4:requests.session
5:~~~~~~~~~~~~~~~~
6:
7:This module provides a Session object to manage and persist settings across
8:requests (cookies, auth, proxies).
9:
10:"""
11:import os
12:from collections import Mapping
13:from datetime import datetime
14:
15:from .auth import _basic_auth_str
16:from .compat import cookielib, OrderedDict, urljoin, urlparse, builtin_str
17:from .cookies import (
18:    cookiejar_from_dict, extract_cookies_to_jar, RequestsCookieJar, merge_cookies)
19:from .models import Request, PreparedRequest, DEFAULT_REDIRECT_LIMIT
20:from .hooks import default_hooks, dispatch_hook
21:from .utils import to_key_val_list, default_headers, to_native_string
22:from .exceptions import (
23:    TooManyRedirects, InvalidSchema, ChunkedEncodingError, ContentDecodingError)
24:from .structures import CaseInsensitiveDict
25:
26:from .adapters import HTTPAdapter
27:
28:from .utils import (
29:    requote_uri, get_environ_proxies, get_netrc_auth, should_bypass_proxies,
30:    get_auth_from_url
31:)
32:
33:from .status_codes import codes
34:
35:# formerly defined here, reexposed here for backward compatibility
36:from .models import REDIRECT_STATI
37:
38:
39:def merge_setting(request_setting, session_setting, dict_class=OrderedDict):
40:    """
41:    Determines appropriate setting for a given request, taking into account the
42:    explicit setting on that request, and the setting in the session. If a
43:    setting is a dictionary, they will be merged together using `dict_class`
44:    """
45:
46:    if session_setting is None:
47:        return request_setting
48:
49:    if request_setting is None:
50:        return session_setting
51:
52:    # Bypass if not a dictionary (e.g. verify)
53:    if not (
54:            isinstance(session_setting, Mapping) and
55:            isinstance(request_setting, Mapping)
56:    ):
57:        return request_setting
58:
59:    merged_setting = dict_class(to_key_val_list(session_setting))
60:    merged_setting.update(to_key_val_list(request_setting))
61:
62:    # Remove keys that are set to None.
63:    for (k, v) in request_setting.items():
64:        if v is None:
65:            del merged_setting[k]
66:
67:    merged_setting = dict((k, v) for (k, v) in merged_setting.items() if v is not None)
68:
69:    return merged_setting
70:
71:
72:def merge_hooks(request_hooks, session_hooks, dict_class=OrderedDict):
73:    """
74:    Properly merges both requests and session hooks.
75:
76:    This is necessary because when request_hooks == {'response': []}, the
77:    merge breaks Session hooks entirely.
78:    """
79:    if session_hooks is None or session_hooks.get('response') == []:
80:        return request_hooks
81:
82:    if request_hooks is None or request_hooks.get('response') == []:
83:        return session_hooks
84:
85:    return merge_setting(request_hooks, session_hooks, dict_class)
86:
87:
88:class SessionRedirectMixin(object):
89:    def resolve_redirects(self, resp, req, stream=False, timeout=None,
90:                          verify=True, cert=None, proxies=None):
91:        """Receives a Response. Returns a generator of Responses."""
92:
93:        i = 0
94:        hist = [] # keep track of history
95:
96:        while resp.is_redirect:
97:            prepared_request = req.copy()
98:
99:            if i > 0:
100:                # Update history and keep track of redirects.
(571 more lines below)

(Open file: n/a)
(Current directory: /psf__requests)
bash-$
\end{Verbatim}
\end{observationbox}
\begin{agentbox}[psf\_\_requests-2317]
The `sessions.py` file is now open, but the `builtin\_str` function call is not visible in the first 100 lines of the file. Since we don't have the exact line number where the `builtin\_str` function is called, we'll need to search within the file for it.



Let's use the `search\_file` command to find the usage of `builtin\_str` within `sessions.py`.

\hrulefill
\begin{CodeVerbatim}[breaklines=true, breakanywhere=true]
search_file "builtin_str"
\end{CodeVerbatim}
\end{agentbox}
\begin{observationbox}[psf\_\_requests-2317]\begin{Verbatim}[breaklines=true, breakanywhere=true]
Found 2 matches for "builtin_str" in /psf__requests/requests/sessions.py:
Line 16:from .compat import cookielib, OrderedDict, urljoin, urlparse, builtin_str
Line 428:        method = builtin_str(method)
End of matches for "builtin_str" in /psf__requests/requests/sessions.py

(Open file: /psf__requests/requests/sessions.py)
(Current directory: /psf__requests)
bash-$
\end{Verbatim}
\end{observationbox}
\begin{agentbox}[psf\_\_requests-2317]
We have found two occurrences of `builtin\_str` in `sessions.py`. The second occurrence on line 428 is where the method is being converted to a string. We should inspect this line to understand the context and see how we can address the issue.



Let's navigate to line 428 in `sessions.py` to see the relevant code.

\hrulefill
\begin{CodeVerbatim}[breaklines=true, breakanywhere=true]
goto 428
\end{CodeVerbatim}
\end{agentbox}
\begin{observationbox}[psf\_\_requests-2317]\begin{Verbatim}[breaklines=true, breakanywhere=true]
[File: /psf__requests/requests/sessions.py (671 lines total)]
(412 more lines above)
413:            data before giving up, as a float, or a (`connect timeout, read
414:            timeout <user/advanced.html#timeouts>`_) tuple.
415:        :type timeout: float or tuple
416:        :param allow_redirects: (optional) Set to True by default.
417:        :type allow_redirects: bool
418:        :param proxies: (optional) Dictionary mapping protocol to the URL of
419:            the proxy.
420:        :param stream: (optional) whether to immediately download the response
421:            content. Defaults to ``False``.
422:        :param verify: (optional) if ``True``, the SSL cert will be verified.
423:            A CA_BUNDLE path can also be provided.
424:        :param cert: (optional) if String, path to ssl client cert file (.pem).
425:            If Tuple, ('cert', 'key') pair.
426:        """
427:
428:        method = builtin_str(method)
429:
430:        # Create the Request.
431:        req = Request(
432:            method = method.upper(),
433:            url = url,
434:            headers = headers,
435:            files = files,
436:            data = data or {},
437:            json = json,
438:            params = params or {},
439:            auth = auth,
440:            cookies = cookies,
441:            hooks = hooks,
442:        )
443:        prep = self.prepare_request(req)
444:
445:        proxies = proxies or {}
446:
447:        settings = self.merge_environment_settings(
448:            prep.url, proxies, stream, verify, cert
449:        )
450:
451:        # Send the request.
452:        send_kwargs = {
453:            'timeout': timeout,
454:            'allow_redirects': allow_redirects,
455:        }
456:        send_kwargs.update(settings)
457:        resp = self.send(prep, **send_kwargs)
458:
459:        return resp
460:
461:    def get(self, url, **kwargs):
462:        """Sends a GET request. Returns :class:`Response` object.
463:
464:        :param url: URL for the new :class:`Request` object.
465:        :param \*\*kwargs: Optional arguments that ``request`` takes.
466:        """
467:
468:        kwargs.setdefault('allow_redirects', True)
469:        return self.request('GET', url, **kwargs)
470:
471:    def options(self, url, **kwargs):
472:        """Sends a OPTIONS request. Returns :class:`Response` object.
473:
474:        :param url: URL for the new :class:`Request` object.
475:        :param \*\*kwargs: Optional arguments that ``request`` takes.
476:        """
477:
478:        kwargs.setdefault('allow_redirects', True)
479:        return self.request('OPTIONS', url, **kwargs)
480:
481:    def head(self, url, **kwargs):
482:        """Sends a HEAD request. Returns :class:`Response` object.
483:
484:        :param url: URL for the new :class:`Request` object.
485:        :param \*\*kwargs: Optional arguments that ``request`` takes.
486:        """
487:
488:        kwargs.setdefault('allow_redirects', False)
489:        return self.request('HEAD', url, **kwargs)
490:
491:    def post(self, url, data=None, json=None, **kwargs):
492:        """Sends a POST request. Returns :class:`Response` object.
493:
494:        :param url: URL for the new :class:`Request` object.
495:        :param data: (optional) Dictionary, bytes, or file-like object to send in the body of the :class:`Request`.
496:        :param json: (optional) json to send in the body of the :class:`Request`.
497:        :param \*\*kwargs: Optional arguments that ``request`` takes.
498:        """
499:
500:        return self.request('POST', url, data=data, json=json, **kwargs)
501:
502:    def put(self, url, data=None, **kwargs):
503:        """Sends a PUT request. Returns :class:`Response` object.
504:
505:        :param url: URL for the new :class:`Request` object.
506:        :param data: (optional) Dictionary, bytes, or file-like object to send in the body of the :class:`Request`.
507:        :param \*\*kwargs: Optional arguments that ``request`` takes.
508:        """
509:
510:        return self.request('PUT', url, data=data, **kwargs)
511:
512:    def patch(self, url, data=None, **kwargs):
(159 more lines below)

(Open file: /psf__requests/requests/sessions.py)
(Current directory: /psf__requests)
bash-$
\end{Verbatim}
\end{observationbox}
\begin{agentbox}[psf\_\_requests-2317]
We are now looking at the section of the `sessions.py` file where the `builtin\_str` function is used to convert the `method` variable. According to the issue, this conversion is causing a problem when the `method` is a binary string because it's being converted to a literal string that includes the binary prefix `b'...'`.



To address this issue, we need to ensure that the `method` variable is correctly converted to a string without the binary prefix. This might involve checking the type of the `method` variable and handling binary strings differently from regular strings.



Let's edit the line to add a check that ensures binary strings are decoded to regular strings before being passed to `builtin\_str`.

\hrulefill
\begin{CodeVerbatim}[breaklines=true, breakanywhere=true]
edit 428:428
        # Ensure method is a string, decode if it's a binary string
        if isinstance(method, bytes):
            method = method.decode('ascii')
        method = builtin_str(method)
end_of_edit
\end{CodeVerbatim}
\end{agentbox}
\begin{observationbox}[psf\_\_requests-2317]\begin{Verbatim}[breaklines=true, breakanywhere=true]
[File: /psf__requests/requests/sessions.py (674 lines total)]
(377 more lines above)
378:    def request(self, method, url,
379:        params=None,
380:        data=None,
381:        headers=None,
382:        cookies=None,
383:        files=None,
384:        auth=None,
385:        timeout=None,
386:        allow_redirects=True,
387:        proxies=None,
388:        hooks=None,
389:        stream=None,
390:        verify=None,
391:        cert=None,
392:        json=None):
393:        """Constructs a :class:`Request <Request>`, prepares it and sends it.
394:        Returns :class:`Response <Response>` object.
395:
396:        :param method: method for the new :class:`Request` object.
397:        :param url: URL for the new :class:`Request` object.
398:        :param params: (optional) Dictionary or bytes to be sent in the query
399:            string for the :class:`Request`.
400:        :param data: (optional) Dictionary or bytes to send in the body of the
401:            :class:`Request`.
402:        :param json: (optional) json to send in the body of the
403:            :class:`Request`.
404:        :param headers: (optional) Dictionary of HTTP Headers to send with the
405:            :class:`Request`.
406:        :param cookies: (optional) Dict or CookieJar object to send with the
407:            :class:`Request`.
408:        :param files: (optional) Dictionary of ``'filename': file-like-objects``
409:            for multipart encoding upload.
410:        :param auth: (optional) Auth tuple or callable to enable
411:            Basic/Digest/Custom HTTP Auth.
412:        :param timeout: (optional) How long to wait for the server to send
413:            data before giving up, as a float, or a (`connect timeout, read
414:            timeout <user/advanced.html#timeouts>`_) tuple.
415:        :type timeout: float or tuple
416:        :param allow_redirects: (optional) Set to True by default.
417:        :type allow_redirects: bool
418:        :param proxies: (optional) Dictionary mapping protocol to the URL of
419:            the proxy.
420:        :param stream: (optional) whether to immediately download the response
421:            content. Defaults to ``False``.
422:        :param verify: (optional) if ``True``, the SSL cert will be verified.
423:            A CA_BUNDLE path can also be provided.
424:        :param cert: (optional) if String, path to ssl client cert file (.pem).
425:            If Tuple, ('cert', 'key') pair.
426:        """
427:
428:        # Ensure method is a string, decode if it's a binary string
429:        if isinstance(method, bytes):
430:            method = method.decode('ascii')
431:        method = builtin_str(method)
432:
433:        # Create the Request.
434:        req = Request(
435:            method = method.upper(),
436:            url = url,
437:            headers = headers,
438:            files = files,
439:            data = data or {},
440:            json = json,
441:            params = params or {},
442:            auth = auth,
443:            cookies = cookies,
444:            hooks = hooks,
445:        )
446:        prep = self.prepare_request(req)
447:
448:        proxies = proxies or {}
449:
450:        settings = self.merge_environment_settings(
451:            prep.url, proxies, stream, verify, cert
452:        )
453:
454:        # Send the request.
455:        send_kwargs = {
456:            'timeout': timeout,
457:            'allow_redirects': allow_redirects,
458:        }
459:        send_kwargs.update(settings)
460:        resp = self.send(prep, **send_kwargs)
461:
462:        return resp
463:
464:    def get(self, url, **kwargs):
465:        """Sends a GET request. Returns :class:`Response` object.
466:
467:        :param url: URL for the new :class:`Request` object.
468:        :param \*\*kwargs: Optional arguments that ``request`` takes.
469:        """
470:
471:        kwargs.setdefault('allow_redirects', True)
472:        return self.request('GET', url, **kwargs)
473:
474:    def options(self, url, **kwargs):
475:        """Sends a OPTIONS request. Returns :class:`Response` object.
476:
477:        :param url: URL for the new :class:`Request` object.
(197 more lines below)
File updated. Please review the changes and make sure they are correct (correct indentation, no duplicate lines, etc). Edit the file again if necessary.

(Open file: /psf__requests/requests/sessions.py)
(Current directory: /psf__requests)
bash-$
\end{Verbatim}
\end{observationbox}
\begin{agentbox}[psf\_\_requests-2317]
The edit has been made to decode binary strings before converting them to a string with `builtin\_str`. This should prevent the issue where a binary string like `b'GET'` is incorrectly converted to the literal string `''b'GET'''`.



To verify that this change resolves the issue, we should create a script that replicates the problem as described. We'll write a script that uses the `requests` library to make a request with a binary method and check if it now handles the method correctly.



Let's create a new file called `reproduce\_issue.py` to test the fix.

\hrulefill
\begin{CodeVerbatim}[breaklines=true, breakanywhere=true]
create reproduce_issue.py
\end{CodeVerbatim}
\end{agentbox}
\begin{observationbox}[psf\_\_requests-2317]\begin{Verbatim}[breaklines=true, breakanywhere=true]
[File: /psf__requests/reproduce_issue.py (1 lines total)]
1:

(Open file: /psf__requests/requests/sessions.py)
(Current directory: /psf__requests)
bash-$
\end{Verbatim}
\end{observationbox}
\begin{agentbox}[psf\_\_requests-2317]
Now that we have created the `reproduce\_issue.py` file, let's write a script that attempts to send a request with a binary method string, which should now be handled correctly due to our fix in `sessions.py`.



We'll write a simple script that imports the `requests` library, sets a binary method, and tries to make a request to a test URL.



\hrulefill
\begin{CodeVerbatim}[breaklines=true, breakanywhere=true]
edit 1:1
import requests

# Define a binary method string
method = b'GET'

# Test URL (example)
url = 'https://httpbin.org/get'

# Attempt to send a request with a binary method
response = requests.request(method, url)

# Print the response status code
print('Response status code:', response.status_code)

# Print a success message if the script completes
print("Script completed successfully, no errors.")
end_of_edit
\end{CodeVerbatim}
\end{agentbox}
\begin{observationbox}[psf\_\_requests-2317]\begin{Verbatim}[breaklines=true, breakanywhere=true]
[File: /psf__requests/reproduce_issue.py (16 lines total)]
1:import requests
2:
3:# Define a binary method string
4:method = b'GET'
5:
6:# Test URL (example)
7:url = 'https://httpbin.org/get'
8:
9:# Attempt to send a request with a binary method
10:response = requests.request(method, url)
11:
12:# Print the response status code
13:print('Response status code:', response.status_code)
14:
15:# Print a success message if the script completes
16:print("Script completed successfully, no errors.")
File updated. Please review the changes and make sure they are correct (correct indentation, no duplicate lines, etc). Edit the file again if necessary.

(Open file: /psf__requests/reproduce_issue.py)
(Current directory: /psf__requests)
bash-$
\end{Verbatim}
\end{observationbox}
\begin{agentbox}[psf\_\_requests-2317]
The script has been written and should now attempt to send a request using a binary method string. This will help us verify if the issue has been resolved by the changes we made to `sessions.py`.



Let's run the script to see if it executes successfully without encountering the issue described.

\hrulefill
\begin{CodeVerbatim}[breaklines=true, breakanywhere=true]
python reproduce_issue.py
\end{CodeVerbatim}
\end{agentbox}
\begin{observationbox}[psf\_\_requests-2317]\begin{Verbatim}[breaklines=true, breakanywhere=true]
Response status code: 200
Script completed successfully, no errors.

(Open file: /psf__requests/reproduce_issue.py)
(Current directory: /psf__requests)
bash-$
\end{Verbatim}
\end{observationbox}
\begin{agentbox}[psf\_\_requests-2317]
The script ran successfully and received a 200 OK response, which indicates that the issue with handling binary method strings has been resolved. The `requests` library is now correctly processing the binary string `b'GET'` as a valid HTTP method.



Since the fix appears to be working, we can clean up by removing the `reproduce\_issue.py` script, as it is no longer needed. After that, we can submit our changes to the code base.



Let's remove the `reproduce\_issue.py` file.

\hrulefill
\begin{CodeVerbatim}[breaklines=true, breakanywhere=true]
rm reproduce_issue.py
\end{CodeVerbatim}
\end{agentbox}
\begin{observationbox}[psf\_\_requests-2317]\begin{Verbatim}[breaklines=true, breakanywhere=true]

(Open file: /psf__requests/reproduce_issue.py)
(Current directory: /psf__requests)
bash-$
\end{Verbatim}
\end{observationbox}
\begin{agentbox}[psf\_\_requests-2317]
The `reproduce\_issue.py` file has been successfully removed. Now that we have verified the fix and cleaned up the test script, we can proceed to submit our changes to the code base.



Let's submit the changes using the `submit` command.

\hrulefill
\begin{CodeVerbatim}[breaklines=true, breakanywhere=true]
submit
\end{CodeVerbatim}
\end{agentbox}
\begin{observationbox}[psf\_\_requests-2317]
\begin{CodeVerbatim}[breaklines=true, breakanywhere=true]
diff --git a/requests/sessions.py b/requests/sessions.py
index d701ff2e..4829747d 100644
--- a/requests/sessions.py
+++ b/requests/sessions.py
@@ -425,6 +425,9 @@ class Session(SessionRedirectMixin):
             If Tuple, ('cert', 'key') pair.
         """
 
+        # Ensure method is a string, decode if it's a binary string
+        if isinstance(method, bytes):
+            method = method.decode('ascii')
         method = builtin_str(method)
 
         # Create the Request.

(Open file: /psf__requests/reproduce_issue.py)
(Current directory: /psf__requests)
bash-\$
\end{CodeVerbatim}
\end{observationbox}

\begin{goldpatchbox}[psf\_\_requests-2317]
\begin{CodeVerbatim}[breaklines=true, breakanywhere=true]
diff --git a/requests/sessions.py b/requests/sessions.py
--- a/requests/sessions.py
+++ b/requests/sessions.py
@@ -13,7 +13,7 @@
 from datetime import datetime
 
 from .auth import _basic_auth_str
-from .compat import cookielib, OrderedDict, urljoin, urlparse, builtin_str
+from .compat import cookielib, OrderedDict, urljoin, urlparse
 from .cookies import (
     cookiejar_from_dict, extract_cookies_to_jar, RequestsCookieJar, merge_cookies)
 from .models import Request, PreparedRequest, DEFAULT_REDIRECT_LIMIT
@@ -425,7 +425,7 @@ def request(self, method, url,
             If Tuple, ('cert', 'key') pair.
         """
 
-        method = builtin_str(method)
+        method = to_native_string(method)
 
         # Create the Request.
         req = Request(
\end{CodeVerbatim}
\end{goldpatchbox}
}
